%!TEX TS-program = pdflatex
%!TEX encoding = UTF-8 Unicode
%!TEX spellcheck = English (Aspell)

\documentclass[11pt]{article}
\usepackage{jeffe,handout,graphicx}

\def\Cdot{\mathbin{\text{\normalfont \textbullet}}}
\def\Sym#1{\textbf{\texttt{\color{BrickRed}#1}}}

\newcommand{\X}[1]{\textcolor{Red}{\mathbf{\underline{#1}}}}

\hidesolutions

% =========================================================
\begin{document}

\headers{CSCI 432/532}{Problem Session 12}{Spring 2026}

\begin{enumerate}\itemsep3ex\parindent1.5em
\item Let's get a sense of how large powers of two really are.
\begin{enumerate}
    \item What is $2^{10}$ exactly?
    \item Assume that $2^{10}=10^3$, one thousand. Using exponent rules, what are $2^{20}$ and $2^{30}$ as powers of 10?
    What rule can you apply to estimate $2^{10k}$ (at least for small positive integers $k$)?
    \item Using the internet, fill in the blanks in the following table. For the last two lines, choose your own
    examples!
   \begin{table}[h]
\centering
\begin{tabular}{|c|l|}
\hline
\textbf{Power of 2} &  \textbf{Interpretation} \\
\hline
$2^{66}$  & Number of cells in a human body \\
\hline
  & Number of different 128-bit RSA keys ($N$) \\
\hline
 & Number of atoms in the observable universe \\
\hline
 & Largest-bit RSA key ever factored \\
\hline
 &  \\
\hline
 &  \\
\hline
\end{tabular}
\end{table} 
\end{enumerate}
\item Assume the time to break an RSA key of length $n$
bits is proportional to $2^{n/2}$ and suppose a 128-bit RSA key takes 1 day to break.
\begin{enumerate}
    \item 
How long would it take to break a 256-bit key under this model?
\item Suppose that, due to faster hardware, computing power doubles every two years.
How long would it take to break a 256-bit key in 20 years?
\item How many bits should you add to the key to maintain the same level of security in 20 years?
\end{enumerate}

\item Now let's assume that the time to break an RSA key of
length $n$ bits is proportional to $e^{n^{1/3}}$ (this is a very rough approximation
of the asymptotic runtime of the Generalized Number Field Sieve algorithm for factoring).
You would again like to answer the question: if computing power doubles every two years,
how many bits should you add to the key to to maintain the same level of security in 20 years?
You may simply evaluate $e^{n^{1/3}}$ for various values of $n$ to get a sense
of how the runtime grows with $n$ and then extrapolate from there.



\item Prove the following:
\begin{enumerate}
    \item If $c|a$ then $c|sa$ for any integer $s$.
    \item If $c|a$ and $c|b$ then $c|a+b$.
    \item If $c|a$ and $c|b$ then $c|sa+tb$ for any integers $s,t$.
    \item If $b>0$, then $\gcd(a,b)=\gcd(a,r)$ where $r$ is the remainder when $a$ is divided by $b$.
\end{enumerate}
\end{enumerate}
\end{document}

